\ssection{BREW-RC}
\label{sec:brew-rc}
We now present our method for bit-exact recovery of ECC in ReRAM crossbars, BREW-RC.
% First, we relate the observed timing behavior—captured as data transitions and their clustering assignments to a latent model defined by an ECC parity matrix, which governs parity-bit switching under a linear block code (LBC).
% Next, we introduce a test-vector generation strategy that guarantees each write-latency measurement ($t_w$) produces a corresponding logical constraint. 
% Finally, we present the resulting SAT-based recovery results. In the following section, we use these results to interpret related work on the MB85AS4MT and MB85AS8MT memories.
% We note here that we must make an assumption that the latent behavior observed in our prior experiments is an ECC function. We admit that this assumption is slightly weaker than the one made in our motivating work \cite{patel2020bit}. In \cite{patel2020bit}, the authors test only modules that advertise ECC capabilities. The datasheet of both devices under test make no mention of error correcting codes, whereas this is common practice for DRAM modules. 
% Our interests, however, reside in the domain of hardware side-channels where an adversary's initial objective is to construct a model of an information leakage channel. Our motivation is therefore related to BEER but not identical.
% \ssubsection{Polarity Relationships and Measurement Pruning}
Section~\ref{sec:clustering} demonstrates that timing measurement distributions identify polarity relationships of entire codewords ($\{\texttt{NONE},\texttt{UNIPOLAR}, \texttt{BIPOLAR}\}$). These distributions are defined with labels:
{
\setlength{\abovedisplayskip}{1pt}
\setlength{\belowdisplayskip}{0.5pt}
\begin{equation}
\ell_{cw}, \ell_d \in \{\texttt{NONE}, \texttt{UNIPOLAR\_P}, \texttt{UNIPOLAR\_N}, \texttt{BIPOLAR}\}.
\end{equation}
} 
where the label ($\ell_{cw}$) reflects the codeword DHD and $\ell_d$ is the known transition polarity of the applied data. Note that a conversion from $\texttt{UNIPOLAR}$ observations to a corresponding $\texttt{UNIPOLAR\_P}$ / $ \texttt{UNIPOLAR\_N}$ label may be required (i.e. 4M, Figure \ref{fig:sierp-hist}). This is trivial as the transition polarity of a unipolar codeword must be identical to the transition polarity of the dataword. We describe the relationship between the applied data and codeword transition label with the shorthand $\ell_{d} \rightarrow\ell_{cw}$.

% Some transition relationships between known dataword DHD and observed codeword DHD derived from $t_w$ measurements are consistent
% (ex. $\texttt{NONE}$ $\rightarrow\texttt{NONE}$, $\texttt{BIPOLAR}\rightarrow\texttt{BIPOLAR}$,  $\texttt{UNIPOLAR}_{N/P}\rightarrow\texttt{UNIPOLAR}_{N/P}$). 
$\texttt{NONE}\rightarrow\texttt{NONE}$ transitions reveal nothing about the parity function. As we showed in Section \ref{sec:red}, the timing of these transitions is related to write serialization overhead and has almost zero variance. 
$\texttt{BIPOLAR}\rightarrow\texttt{BIPOLAR}$ also gives no information. These observations essentially translate to: ``given a forced \texttt{BIPOLAR} multi-bit data transition, the parity bits may or may not have transitioned in any polarity.''
Some observations would violate our hypothesis that parity bit polarity serialization is the cause of our observed timing fluctuations (ex. $\texttt{NONE} \to \texttt{UNIPOLAR}_{N/P}$,  $\texttt{BIPOLAR}\rightarrow\texttt{UNIPOLAR}_{N/P}$ ). We implement assertions to validate that they are never observed.

BREW-RC uses two observation classes to generate SAT constraints.
$\texttt{UNIPOLAR}_{N/P}\rightarrow\texttt{UNIPOLAR}_{N/P}$ reveal an explicit constraint: ``given a $\texttt{UNIPOLAR}_{N/P}$, data transition, no parity bits transition in an opposing polarity.'' $\texttt{UNIPOLAR}_{N/P}\rightarrow\texttt{BIPOLAR}$ produce an existential constraint that translates to: ``given a $\texttt{UNIPOLAR}_{N/P}$ data transition, at least one parity bit transitioned in an opposing polarity.''

\ssubsection{Modeling parity-bit activity}
As described in Section~\ref{sec:bg}, ReRAM memories often employ linear block codes (LBCs). For a dataword $d\in\{0,1\}^k$, the codeword $c=dG=[d\mid p]$ propagates data bits ($d$) and appends generated parity bits generated by the $GF(2)$ multiplication
{
\setlength{\abovedisplayskip}{1pt}
\setlength{\belowdisplayskip}{1pt}
\begin{equation}
    p = dP.
\end{equation}
}
For pre- and post-write states $(d_0,d_1)$, the parity-bit updates are
{
\setlength{\abovedisplayskip}{1pt}
\setlength{\belowdisplayskip}{1pt}
\begin{equation}
    \Delta p = (d_0 + d_1)P = \Delta d P,
\end{equation}
}
where $\Delta d=d_0+d_1$ marks toggled data bits.  
Thus, $\Delta p$ acts as a logical gate that determines which parity bits update under a given data transition, independent of transition polarity.  
In the SAT construction, $\Delta p$ is evaluated for every dataword transition and used in all clauses to restrict parity-bit activity to the toggled subset.

\ssubsection{Polarity-encoded SAT formulation}
\label{sec:sat-encoding}

To incorporate polarity, each transition is assigned a data-polarity label $\ell_d \in \{\texttt{UNIPOLAR\_P},\texttt{UNIPOLAR\_N}\}$ 
with scalar representation $d\in\{-1,+1\}$:
$d=+1$ for positive ($0{\rightarrow}1$) and $d=-1$ for negative ($1{\rightarrow}0$) unipolar transitions.
Bipolar and static cases (\texttt{BIPOLAR}, \texttt{NONE}) produce no directional information and are excluded.

For each parity bit $p_j$, Boolean variables $p_{0,j}$ and $p_{1,j}$ denote its pre- and post-write values, gated by $\Delta p_j$. 
We then define
{
\setlength{\abovedisplayskip}{1pt}
\setlength{\belowdisplayskip}{1pt}
\begin{align}
    \textit{pos}_j &\leftrightarrow \Delta p_j \wedge (\neg p_{0,j} \wedge p_{1,j}), \\
    \textit{neg}_j &\leftrightarrow \Delta p_j \wedge (p_{0,j} \wedge \neg p_{1,j}).
\end{align}
}
encoding positive and negative parity transitions only when the parity bit is expected to update.  
Depending on $(\ell_d,\ell_{cw})$, one of four clause types is generated as summarized in Table~\ref{tab:polarity-clauses}.  
Under unipolar test patterns, the \texttt{NONE} cluster is unobservable (more than zero bits are active). 
A \texttt{BIPOLAR} observation indicates that at least one parity bit transitioned in the opposite polarity to the data.
Conversely, a matching unipolar observation indicates the absence of any opposite-polarity parity transition.

Each $(d_0,d_1,\ell_d,\ell_{cw})$ tuple contributes a compact SAT subproblem constraining the allowed polarity of parity-bit transitions under a candidate parity submatrix ($P$).  

Because prior measurements exercised only single-bit or single-byte toggles, the resulting clauses may not reveal cross-byte parity coupling. Additionally, redundant writes and bipolar dataword transitions that are not useful, as they do not emit SAT constraints. Therefore, we introduce a \emph{unipolar random word walk} that systematically explores unipolar transitions over word regions, efficiently sampling the full word-level unipolar transition space.
\begin{table}[t!]
\setlength{\textfloatsep}{0pt}
\setlength{\floatsep}{0pt}
\centering
\caption{SAT solver clause definitions \label{tab:polarity-clauses}.}
\vspace{-.2cm}
\renewcommand{\arraystretch}{1.1}
\setlength{\tabcolsep}{3pt}
\begin{tabular*}{\linewidth}{@{\extracolsep{\fill}}lcccc@{}}
\toprule
\textbf{Data label} $\ell_d$ & \textbf{Obs.} $\ell_{cw}$ & \textbf{Constraint} & \textbf{Description} \\
\midrule
\texttt{UNIPOLAR\_P} & \texttt{UNIPOLAR\_P} & $\forall j\, (p_{0,j}\vee\neg p_{1,j})$  & forbid \emph{neg}  \\
\texttt{UNIPOLAR\_N} & \texttt{UNIPOLAR\_N} & $\forall j\, (\neg p_{0,j}\vee p_{1,j})$ & forbid \emph{pos} \\
\texttt{UNIPOLAR\_P} & \texttt{BIPOLAR}     & $\bigvee_j \textit{neg}_j$               & require $\ge1$ \emph{neg}  \\
\texttt{UNIPOLAR\_N} & \texttt{BIPOLAR}     & $\bigvee_j \textit{pos}_j$               & require $\ge1$ \emph{pos}  \\
\bottomrule
\end{tabular*}  
\end{table}
\setlength{\textfloatsep}{-10pt} {
	% \setlength{\floatsep}{-10pt}
	\begin{algorithm}[t!]
		\footnotesize
		\caption{1-word, random, unipolar walk\label{alg:ruw}}
		\begin{algorithmic}[1]
			\Require word size $k$, length $\ell$, max test-vectors per address $\textit{max\_tv}$
			\State $\texttt{Keys} \gets \varnothing$; $\texttt{A} \gets \varnothing$; \Comment{used transitions, used addrs.}
			\While{$i < \ell$}
			\If{$i \bmod \text{max\_tv} = 0$}
			\State $\texttt{addr}\gets\texttt{GetRandWordAddr(A)}$; $\texttt{A} \cup \texttt{addr}$;
			\State $s \gets \texttt{GetRandWord()}$;
			\State write\_word(addr, $s$)
			\EndIf
			\State $d \gets  s=0 \; ? \; 1 : s=2^k-1 \; ? \; 0 : \texttt{RNG(0,1)}))$;
			\State $m \gets \texttt{RNG(1, $2^k -1$)}$
			\State $s' \gets d=0\; ? \;s\; \& \; \neg m : s\; | \; m $;
			\State $key \gets \min(s,s') << k \; | \; max(s, s')$;
			\If{$\texttt{key} \notin \texttt{Keys}$}
			\State $\texttt{Keys} \cup \texttt{k}$ ; $s \gets s'$ ; $i\gets i+1$;
			\State write\_word($\texttt{addr}$, $s'$)
			\EndIf
			\EndWhile
		\end{algorithmic}
	\end{algorithm}
}
\setlength{\textfloatsep}{3pt}

\ssubsection{BREW-RC Data Collection}
\label{exp:2ii}
To ensure every test vector emits a useful constraint, we perform a random unipolar walk of the dataword space, and measure $t_{w}$ for every applied test vector. The details of this test-vector generation technique are shown in Alg. \ref{alg:ruw}. This procedure ensures that a unique unipolar dataword transition is generated for every write operation in the programming sequence.

% A random word-aligned address in memory is selected, and initialized to a random state. Then a direction (0 for \textit{neg}, 1 for \textit{pos}) is randomly selected. A random mask is generated according to the polarity, and then applied to the last test vector to generate the next unipolar random transition. This process repeats until the desired number of transitions have been applied. To ensure selected addresses are not worn out over the course of long experiments, we log wear levels and select new addresses once a threshold has been met. 

The 4M has a much higher probability of generating $\texttt{UNIPOLAR}_{P/N}$ $\to \texttt{UNIPOLAR}_{P/N}$ constraints due to it's significantly smaller parity submatrix than the 8M. For each chip, we repeat the experiment 10 times at different word addresses to generate a representative $t_{wmin}$. In order for BREW-RC to converge to a single ECC, we collect 100K (200K) $t_{wmin}$ from each 4M (8M) sample. 

After capturing timing data, we assign labels to construct $(d_0,d_1,$ $\ell_d,$ $\ell_{cw}$) observations. Each observation yields a set of XOR constraints encoding $\Delta p$ and CNF clauses enforcing the allowed polarity of each parity-bit transition. The resulting hybrid XOR+CNF formula is solved using CryptoMiniSAT \cite{soos2009extending, soos2018cryptominisat5}.

% Figure \ref{fig:1w-rand-walk} shows the distribution of $\min\left(t_{w}\right)$ measurements normalized for each chip along with the empirical cumulative distribution function. The upper plots show labels applied by clustering the $\min\left(t_{ws}\right)$ observations ($\ell_{cw}$), and the lower plots show labels applied using only the known data transitions ($\ell_d$). Clustering is performed using a mixture of extreme-value (Gumbel) distributions to capture the timing model discussed in \ref{sec:bg}.
% \begin{figure}[t!]
    \centering
    \includegraphics[width=0.99\columnwidth]{fig/images/single-word-rand-walk-clustering.pdf}
    \Description{Histograms for the 4Mb and 8Mb write latencies measured for the random walk experiment are shown. The X-axis is the mix-max normalized write latency and the y-axis is the log-probability that a write latency is observed. Write latencies appear in three distinct distributions. Each distribution appears as a light-tailed left-skewed distribution. Histograms for \texttt{UNIPOLAR\_N}, \texttt{UNIPOLAR\_P} test vectors are overlaied and shown in different colors. The leftmost (slowest) distribution contains only \texttt{UNIPOLAR\_N}, the middle distribution contains mostly \texttt{UNIPOLAR\_P} (In the 4Mb it overlaps with  \texttt{UNIPOLAR\_N}), and the rightmost distribution for both memories contains a mixture of the two.}
    \caption{Per-chip min-max normalized $\min \left(t_{w}\right)$ measurement histograms and ECDFs for a 1-word random unipolar walk experiment used for BREW-RC constraint generation. Clustering is performed per-chip to demonstrate reproducibility across samples. Labels are assigned either naively from applied test vector data (bottom), or from clustering performed on observed write latencies (top). \label{fig:1w-rand-walk}}
\end{figure}
% To adequately exercise the ECC function, such that we can recover a single ECC function up to column permutations, we require test vectors that exercise all possible intra-bit coupling. We implement a simple unipolar random walk test vector generator shown in Alg. \ref{alg:rand-walk}. The unipolar random-walk test vectors require a known message length ($k$). Therefore BREW-RC requires only two test pattern sequences to fully extract $P$ (Alg. \ref{alg:red} and Alg. \ref{alg:rand-walk}).